% pre-ec2.tex
\documentclass[11pt, a4paper]{article}

\usepackage[T1]{fontenc}
\usepackage[utf8]{inputenc}
\usepackage{tgpagella}
\usepackage{microtype}
\usepackage[spanish, es-noshorthands]{babel}
\usepackage{amsmath, amssymb, bm}
\usepackage[most]{tcolorbox}
\usepackage[margin=2.5cm]{geometry}
\usepackage{fancyhdr}

\pagestyle{fancy}
\fancyhf{}
\setlength{\headheight}{16pt}
\lhead{\textbf{IMT-342}}
\chead{Práctica Integradora Pre-EC2}
\rhead{Independiente}

\definecolor{ucbblue}{RGB}{0, 51, 102}

\begin{document}

\begin{center}
    \Large\textbf{\textcolor{ucbblue}{Especificación de Ejercicio Integrador: Simulación de EC2}}\\[0.2cm]
\end{center}

\begin{tcolorbox}[colback=white, colframe=black, title=\textbf{Objetivo Evaluativo[cite: 14]}]
    \small Esta práctica verifica si usted puede ejecutar de manera \textbf{independiente} los análisis matemáticos clave antes de someterse al examen EC2. Coleccione su evidencia individual antes de debatir con sus compañeros.
\end{tcolorbox}

\section*{A. IK y Pieper (Desacoplamiento)[cite: 14]}
Dada una pose deseada $T_d = \begin{bmatrix} R_d & p_d \\ 0 & 1 \end{bmatrix}$, con $p_d = [0.8, 0, 0.4]^T$ y orientación tal que $z_6^d = [1, 0, 0]^T$. Para un offset final de $d_6 = 0.1$:
\begin{enumerate}
    \item Calcule el centro de muñeca $p_W$.
    \item Si una resolución parcial del brazo nos lleva a $c_2 = 0.5 \implies s_2 = \pm\sqrt{0.75}$. Construya la expresión rigurosa de recuperación de ángulo utilizando $\operatorname{atan2}$.
\end{enumerate}

\section*{B. Jacobiano[cite: 14]}
Evalúe la siguiente construcción hipotética en $q = [0, \pi/2]^T$. Ejes $Z$ coinciden con el vector global $Z_0$:
\begin{equation*}
    J = \begin{bmatrix} \dots & -0.5 \\ \dots & 0 \\ 0 & 0 \\ 0 & 0 \\ 0 & 0 \\ 1 & 1 \end{bmatrix}, \quad \text{si se requiere } \dot{q} = \begin{bmatrix} 0 \\ 2 \end{bmatrix} \implies \text{Calcule } \mathbf{v}_e.
\end{equation*}

\section*{C. Singularidades[cite: 14]}
Para un sub-jacobiano $J_p = \begin{bmatrix} -2s_{12}-s_1 & -2s_{12} \\ 2c_{12}+c_1 & 2c_{12} \end{bmatrix}$. 
Demuestre analíticamente si la configuración $q_2 = \pi$ produce una pérdida de rango. Identifique la dirección de movimiento cartesiano que se pierde.

\section*{D. Trayectoria LSPB[cite: 14]}
Debe moverse de $q_i=0$ a $q_f=4$ en $t_f=4$ s. ¿Es factible configurar un perfil trapezoidal con $\ddot{q}_c=1$ rad/s$^2$? De ser así, determine el tiempo de mezcla de aceleración $t_c$ e interprete las fases del movimiento.

\section*{E. Dinámica de 1R[cite: 14]}
Evaluando un eslabón 1R en el momento pico de la fase de aceleración calculada en LSPB, considerando $J_{\text{eq}}=0.5$, $\ell_c=0.5, m=2$. Calcule el torque requerido en la configuración instantánea de $q=0.5$ rad. Identifique la contribución inercial y la gravitacional.

\section*{F. Razonamiento Físico Integrado[cite: 14]}
Dado el estado instantáneo calculado en la sección E, justifique simultáneamente: ¿Por qué hay alto riesgo dinámico pero bajo riesgo de singularidad en ese punto de la trayectoria?

\end{document}
